% !TeX root = ../BTechProjectTemplate.tex
\chapter{Results \& Discussion}

\section{Quantitative Evaluation}
The performance of Swin-MMHCA is evaluated using Peak Signal-to-Noise Ratio (PSNR) and Structural Similarity Index (SSIM). We compare our framework against several state-of-the-art methods on the IXI dataset for the T2-weighted target modality.

\begin{table}[H]
\centering
\caption{Comparison of Super-Resolution Performance on IXI Dataset.}
\label{tab:quantitative_results}
\resizebox{\columnwidth}{!}{
\begin{tabular}{l|cc|cc}
\toprule
\textbf{Method} & \multicolumn{2}{c|}{\textbf{2$\times$}} & \multicolumn{2}{c}{\textbf{4$\times$}} \\
& PSNR $\uparrow$ & SSIM $\uparrow$ & PSNR $\uparrow$ & SSIM $\uparrow$ \\
\midrule
Bicubic & 33.44 & 0.9589 & 27.86 & 0.8611 \\
SRCNN \cite{srcnn} & 37.32 & 0.9796 & 29.69 & 0.9052 \\
VDSR \cite{vdsr} & 38.65 & 0.9836 & 30.79 & 0.9240 \\
IDN \cite{idn} & 39.09 & 0.9846 & 31.37 & 0.9312 \\
RDN \cite{rdn} & 38.75 & 0.9838 & 31.45 & 0.9324 \\
FSCWRN \cite{fscwrn} & 39.44 & 0.9855 & 31.71 & 0.9359 \\
CSN \cite{csn} & 39.71 & 0.9863 & 32.05 & 0.9413 \\
EDSR \cite{edsr} & 39.81 & 0.9865 & 31.83 & 0.9377 \\
\textbf{Swin-MMHCA (Ours)} & \textbf{40.62} & \textbf{0.9927} & \textbf{33.32} & \textbf{0.9572} \\
\bottomrule
\end{tabular}
}
\end{table}

Our Swin-MMHCA model achieves a superior PSNR and SSIM across both scales, demonstrating the success of the hierarchical transformer architecture in medical SR.

\section{Ablation Study and Modular Performance Analysis}
To evaluate the contribution of each module within the Swin-MMHCA framework, we conducted a rigorous ablation study. We started with a baseline SwinV2 model and progressively added the proposed modules.

\begin{itemize}
    \item \textbf{Baseline SwinV2 (Unimodal):} When trained only on T2-weighted images, the model achieved a PSNR of 38.12 dB for 2$\times$ upscaling and 31.14 dB for 4$\times$. While the transformer captured global context, it lacked the high-frequency structural details often present in other modalities.
    \item \textbf{Multi-modal Concatenation (Early Fusion):} Simply concatenating T1 and PD modalities at the input improved the PSNR to 38.65 dB (+0.53 dB) for 2$\times$ and 31.52 dB (+0.38 dB) for 4$\times$. However, the model struggled to reconcile different intensity distributions.
    \item \textbf{MMHCA Module:} The introduction of the Multi-Modal Hybrid Cross-Attention module was the most significant factor, reaching 40.48 dB (+1.83 dB gain) for 2$\times$ and 32.90 dB (+1.38 dB gain) for 4$\times$. It effectively recovered sharp cortical boundaries by querying relevant features across scans.
    \item \textbf{Edge Guidance Branch:} Adding the Sobel-based edge extractor further refined the structural boundaries, bringing the PSNR to 40.55 dB for 2$\times$ and 33.11 dB for 4$\times$. This module acted as a structural regularizer, ensuring pixel alignment with anatomical gradients.
    \item \textbf{Multi-Task Semantic Supervision:} The final addition of segmentation and lesion detection heads provided a gain to 40.62 dB for 2$\times$ and 33.32 dB for 4$\times$. More importantly, it significantly reduced the LPIPS score, indicating enhanced perceptual quality and anatomical fidelity.
\end{itemize}

\section{Qualitative Analysis and Visual Results}
Visual inspection of the SR outputs across the test set confirms the quantitative findings. Swin-MMHCA demonstrates a superior ability to reconstruct the cerebral cortex's intricate sulci and gyri patterns. We provide visual evidence for both $2\times$ and $4\times$ upscaling factors.

\begin{figure}[htbp]
    \centering
    \includegraphics[width=\textwidth]{figures/comparison_figure_x4.png}
    \caption{Qualitative comparison of $4\times$ super-resolution results on the IXI T2-weighted brain MRI. From left to right: Low-Resolution (4x), Bicubic interpolation, EDSR-Nav (MHCA-main baseline), Swin-MMHCA (Ours), and High-Resolution Ground Truth. The zoom-in boxes highlight the recovery of sharp gray-white matter boundaries and fine anatomical textures.}
    \label{fig:results_visual_x4}
\end{figure}

\begin{figure}[htbp]
    \centering
    \includegraphics[width=\textwidth]{figures/comparison_figure_x2.png}
    \caption{Qualitative comparison of $2\times$ super-resolution results. At this lower upscaling factor, Swin-MMHCA achieves nearly perfect anatomical reconstruction, effectively bridging the gap between clinical-grade and high-resolution research scans.}
    \label{fig:results_visual_x2}
\end{figure}

As shown in Fig. \ref{fig:results_visual_x4}, Bicubic interpolation suffers from significant blurring, making it impossible to delineate the ventricular boundaries accurately. The EDSR-Nav model introduces "ringing" artifacts near high-contrast edges. In contrast, our Swin-MMHCA model produces a result that is visually indistinguishable from the ground truth for most subjects. The cross-attention mechanism successfully "borrows" the sharp contrast from the T1 modality to reconstruct the fine-grained details of the T2 scan, even at a high upscaling factor of $4\times$.

\section{Regional Performance Analysis}
To understand the model's strengths, we analyzed its performance across different anatomical regions of the brain.

\subsection{Cortical Sulci and Gyri}
Swin-MMHCA showed the most significant improvement over baseline methods in the cerebral cortex. The intricate folding patterns of the sulci are often blurred into a uniform gray mass by bicubic interpolation. Our model, guided by the T1 structural reference, successfully reconstructed the sharp boundaries between the gray matter and the cerebrospinal fluid in the subarachnoid space.

\subsection{Ventricular Boundaries}
The lateral ventricles provide high-contrast boundaries that are essential for tracking brain atrophy in neurodegenerative diseases. Swin-MMHCA achieved nearly perfect reconstruction of these fluid-filled spaces, avoiding the "staircase" artifacts that often plague pure CNN architectures.

\section{Clinical Implications of High-Fidelity SR}
The quantitative gains achieved by Swin-MMHCA (e.g., reaching 40.62 dB PSNR) have direct clinical implications. In radiomics, feature extraction algorithms are highly sensitive to image noise and resolution. By producing super-resolved images that are structurally consistent with high-resolution ground truth, our framework enables:
\begin{enumerate}
    \item \textbf{More Accurate Volumetry:} Reducing the partial volume effect leads to more precise measurements of hippocampal volume, a key biomarker for Alzheimer's disease.
    \item \textbf{Improved Lesion Detection:} Sharp textures allow for the detection of subtle micro-lesions in Multiple Sclerosis that might be missed in low-resolution clinical scans.
    \item \textbf{Surgical Planning:} High-resolution maps of the cortical surface are essential for neurosurgical navigation and avoiding critical functional areas.
\end{enumerate}
Our results suggest that Swin-MMHCA can effectively bridge the gap between "fast clinical protocols" and "high-resolution research standards," potentially reducing the need for long, expensive scans.
